% !TEX root = ../algebra_02.tex

\section{Циклические подпространства}

Пусть $\dim V = n < \infty$ и $A \in \operatorname{End} V$. Для любого вектора $v \in V$ введём понятие циклического подпространства.

\begin{Definition}{Циклическое подпространство}{}
	Циклическим подпространством $C_v$ называется минимальное $A$-инвариантное подпространство, содержащее вектор $v$. Эквивалентно, оно определяется как пересечение всех таких подпространств:
	\[
	C_v = \bigcap\{W \subseteq V \mid W \text{ — $A$-инвариантное подпространство и } v \in W\}.
	\]
\end{Definition}

\subsubsection*{Свойство инвариантности}

Поскольку пересечение $A$-инвариантных подпространств снова $A$-инвариантно, подпространство $C_v$ также является $A$-инвариантным. Кроме того, если $W$ — любое $A$-инвариантное подпространство и $v \in W$, то $C_v \subseteq W$.

\subsection{Подстановка оператора в многочлен}

В теории линейных операторов удобно рассматривать многочлены от оператора. Пусть $f \in K[x]$ — многочлен вида $f(x) = a_0 + a_1x + \dots + a_mx^m.$
Тогда оператор $f(A) \in \operatorname{End} V$ определяется следующим образом: $f(A) = a_0 E_V + a_1 A + a_2 A^2 + \dots + a_m A^m,$
где под степенью понимается композиция: $A^k = \underbrace{A \circ A \circ \dots \circ A}_{k}.$

\subsubsection*{Гомоморфизм подстановки}

Отображение $\mathcal{E}_A \colon K[x] \to \operatorname{End} V, \qquad f \mapsto f(A),$
является гомоморфизмом колец, то есть:
\begin{enumerate}
	\item $\mathcal{E}_A(1) = E_V$;
	\item $\mathcal{E}_A(f + g) = \mathcal{E}_A(f) + \mathcal{E}_A(g)$;
	\item $\mathcal{E}_A(fg) = \mathcal{E}_A(f)\mathcal{E}_A(g)$.
\end{enumerate}

\subsection{Аннулятор вектора и базис циклического подпространства}

Пусть $v \in V$. Рассмотрим множество многочленов, которые обнуляют данный вектор после подстановки в них оператора $A$.

\begin{Definition}{Аннулятор вектора}{}
	Множество $I_v = \{ f \in K[x] \mid f(A)v = 0 \}$
	называется аннулятором вектора $v$.
	
	Заметим, что $I_v = \ker \operatorname{ev}_v$, где $\operatorname{ev}_v \colon K[x] \to V, \qquad f \mapsto f(A)v.$
\end{Definition}

\begin{Lemma}{Свойства аннулятора}{}
	$I_v$ является ненулевым идеалом в кольце $K[x]$.
\end{Lemma}

\subsubsection*{Доказательство}

\begin{enumerate}
	\item Если $f, g \in I_v$, то $(f + g)(A)v = f(A)v + g(A)v = 0.$
	Значит, $f + g \in I_v$.
	
	\item Если $g \in I_v$ и $f \in K[x]$, то $(fg)(A)v = f(A)\bigl(g(A)v\bigr) = f(A)0 = 0.$
	Значит, $fg \in I_v$.
	
	\item Покажем, что $I_v \neq \{0\}$. Рассмотрим систему векторов $v, Av, A^2v, \dots, A^nv.$
	Так как $\dim V = n$, эти $n+1$ векторов линейно зависимы. Следовательно, существуют коэффициенты $\alpha_0,\dots,\alpha_n$, не все равные нулю, такие что $\alpha_0 v + \alpha_1 Av + \dots + \alpha_n A^nv = 0.$
	Тогда ненулевой многочлен $f(x) = \alpha_0 + \alpha_1x + \dots + \alpha_nx^n$
	принадлежит $I_v$.
\end{enumerate}

Так как $K[x]$ — область главных идеалов, идеал $I_v$ порождается одним элементом: $I_v = (M_{v,A}).$

\begin{Definition}{Минимальный аннулятор вектора}{}
	Нормированный многочлен минимальной степени из $I_v$, $M_{v,A}(x) = x^d + \beta_{d-1}x^{d-1} + \dots + \beta_0,$
	называется минимальным аннулятором вектора $v$ относительно оператора $A$.
\end{Definition}

\subsection{Базис циклического подпространства}

\begin{Proposition}{}{}
	Пусть $\deg M_{v,A} = d$. Тогда система векторов $v, Av, A^2v, \dots, A^{d-1}v$
	является базисом циклического подпространства $C_v$.
\end{Proposition}

\subsubsection*{Доказательство}

Линейная независимость этой системы следует из минимальности степени многочлена $M_{v,A}$: иначе существовал бы ненулевой многочлен степени меньше $d$, аннулирующий $v$.

Кроме того, $M_{v,A}(A)v = A^dv + \beta_{d-1}A^{d-1}v + \dots + \beta_0v = 0,$
поэтому $A^dv = -\beta_{d-1}A^{d-1}v - \dots - \beta_0v.$
Следовательно, все последующие векторы $A^kv$ выражаются через $v, Av, \dots, A^{d-1}v.$
Значит, эта система порождает $C_v$ и является его базисом.

\subsection{Характеристический многочлен и теорема Гамильтона--Кэли}

Вычислим характеристический многочлен оператора, ограниченного на циклическое подпространство $C_v$. Матрица оператора $A|_{C_v}$ в базисе $v, Av, \dots, A^{d-1}v$
является сопровождающей матрицей для многочлена $M_{v,A}$:
\[
[A|_{C_v}] =
\begin{pmatrix}
	0 & 0 & \dots & 0 & -\beta_0 \\
	1 & 0 & \dots & 0 & -\beta_1 \\
	0 & 1 & \dots & 0 & -\beta_2 \\
	\vdots & \vdots & \ddots & \vdots & \vdots \\
	0 & 0 & \dots & 1 & -\beta_{d-1}
\end{pmatrix}.
\]

Будем использовать стандартное соглашение $\chi_A(x) = \det(xE - A).$
Тогда
\[
\chi_{A|_{C_v}}(x)
=
\det(xE - [A|_{C_v}])
=
x^d + \beta_{d-1}x^{d-1} + \dots + \beta_0
=
M_{v,A}(x).
\]

\begin{Theorem}{Гамильтона--Кэли}{}
	Пусть $\dim V < \infty$ и $A \in \operatorname{End} V$. Тогда $\chi_A(A) = 0.$
\end{Theorem}

\subsubsection*{Идея доказательства}

Для любого вектора $v \in V$:
\begin{enumerate}
	\item Характеристический многочлен оператора на инвариантном подпространстве делит характеристический многочлен всего оператора: $\chi_{A|_{C_v}} \mid \chi_A.$
	
	\item Следовательно, $\chi_A = \chi_{A|_{C_v}} \cdot g = M_{v,A}\cdot g$
	для некоторого $g \in K[x]$.
	
	\item Тогда $\chi_A(A)v = g(A)M_{v,A}(A)v = 0,$
	поскольку $M_{v,A}$ аннулирует вектор $v$.
	
	\item Значит, $\chi_A(A)v = 0$ для любого $v \in V$, откуда $\chi_A(A) = \mathbf{0}.$
\end{enumerate}

\subsection{Минимальный многочлен и разложение пространства}

Из теоремы Гамильтона--Кэли следует, что множество многочленов, аннулирующих оператор $A$, непусто.

\begin{Definition}{Минимальный многочлен оператора}{}
	Минимальным многочленом оператора $A$ называется нормированный многочлен $\mathcal{M}_A$ минимальной степени такой, что $\mathcal{M}_A(A) = \mathbf{0}.$
	Он порождает идеал аннуляторов оператора: $(\mathcal{M}_A) = \{ f \in K[x] \mid f(A) = 0 \}.$
\end{Definition}

\subsubsection*{Свойства}

\begin{enumerate}
	\item Из теоремы Гамильтона--Кэли следует, что $\mathcal{M}_A \mid \chi_A,$
	то есть минимальный многочлен делит характеристический.
	
	\item Характеристический и минимальный многочлены имеют одни и те же корни над алгебраическим замыканием поля $K$, но кратности этих корней могут различаться.
\end{enumerate}

\subsection{Теорема о разложении: лемма о ядрах}

Пусть $(fg)(A) = 0$ и многочлены $f$ и $g$ взаимно просты: $(f,g) = 1.$

\begin{Proposition}{}{}
	В этих условиях пространство $V$ разлагается в прямую сумму $A$-инвариантных подпространств: $V = V_1 \oplus V_2, \qquad V_1 = \ker f(A), \qquad V_2 = \ker g(A).$
	При этом $f(A)|_{V_1} = 0, \qquad g(A)|_{V_2} = 0.$
\end{Proposition}

\subsubsection*{Доказательство}

Так как $(f,g) = 1$, по теореме Безу существуют такие $a,b \in K[x]$, что $fa + gb = 1.$
Подставим оператор $A$: $f(A)a(A) + g(A)b(A) = E.$
Для любого $v \in V$ получаем $v = \underbrace{g(A)b(A)v}_{v_1} + \underbrace{f(A)a(A)v}_{v_2}.$
Здесь $v_1 \in \ker f(A)$, так как $f(A)v_1 = f(A)g(A)b(A)v = (fg)(A)b(A)v = 0.$
Аналогично $v_2 \in \ker g(A)$. Следовательно, $V = \ker f(A) + \ker g(A).$

Осталось проверить, что сумма прямая. Если $v \in \ker f(A) \cap \ker g(A),$
то $v = a(A)f(A)v + b(A)g(A)v = 0.$
Значит, $V = \ker f(A) \oplus \ker g(A).$

\subsection{Обобщение и циклическая структура}

В общем случае, если $\mathcal{M}_A = p_1^{d_1} \dots p_t^{d_t},$
где $p_1,\dots,p_t$ — попарно различные неприводимые многочлены, то $V = W_1 \oplus W_2 \oplus \dots \oplus W_t, \qquad W_j = \ker\bigl(p_j^{d_j}(A)\bigr).$
Каждое $W_j$ в свою очередь разлагается в прямую сумму циклических подпространств: $W_j = W_{j_1} \oplus \dots \oplus W_{j_e}, \qquad W_{j_i} = C_{v_{j_i}}.$

\subsection{Жорданова форма и сопровождающая матрица}

Если неприводимый многочлен линейный, то есть $p_j = x - \lambda_j,$
структуру оператора на соответствующем подпространстве можно свести к изучению нильпотентного оператора $N = A - \lambda_j E.$

В базисе $E = (N^{d-1}v, N^{d-2}v, \dots, Nv, v)$
матрица оператора $A$ на циклическом подпространстве имеет вид жордановой клетки:
\[
[A|_{C_v}]_E =
\begin{pmatrix}
	\lambda_j & 1 & 0 & \dots & 0 \\
	0 & \lambda_j & 1 & \dots & 0 \\
	\vdots & \vdots & \ddots & \ddots & \vdots \\
	0 & 0 & \dots & \lambda_j & 1 \\
	0 & 0 & \dots & 0 & \lambda_j
\end{pmatrix}
= J_d(\lambda_j).
\]

В общем случае матрица оператора будет блочно-диагональной, где блоки соответствуют сопровождающим матрицам для многочленов $f_j$:
\[
L =
\begin{pmatrix}
	L_{j_1} & & & 0 \\
	& L_{j_2} & & \\
	& & \ddots & \\
	0 & & & L_{j_k}
\end{pmatrix}.
\]

Для многочлена $f(x) = x^n + a_{n-1}x^{n-1} + \dots + a_0$
сопровождающая матрица $L_f$ имеет вид
\[
L_f =
\begin{pmatrix}
	0 & 0 & \dots & 0 & -a_0 \\
	1 & 0 & \dots & 0 & -a_1 \\
	0 & 1 & \dots & 0 & -a_2 \\
	\vdots & \vdots & \ddots & \vdots & \vdots \\
	0 & 0 & \dots & 1 & -a_{n-1}
\end{pmatrix}.
\]

\subsection{Теория групп. Простейшие свойства}

Начинаем изучение алгебраических структур с одним законом композиции.

\subsection{Основные свойства операций в группе}

\begin{Proposition}{Закон сокращения}{}
	Пусть $G$ — группа. Тогда для любых элементов $g,g_1,g_2 \in G$ выполняются следующие свойства:
	\begin{enumerate}
		\item $g_1g = g_2g \Rightarrow g_1 = g_2$ \quad (правое сокращение);
		\item $gg_1 = gg_2 \Rightarrow g_1 = g_2$ \quad (левое сокращение).
	\end{enumerate}
\end{Proposition}

\begin{Proposition}{Обобщённая ассоциативность}{}
	Результат умножения в выражении $g_1g_2\dots g_n$
	не зависит от расстановки скобок, если порядок множителей не меняется.
\end{Proposition}

\subsubsection*{Формальная запись}

Произвольное произведение $g_1g_2\dots g_n$ всегда можно привести к левонормированному виду: $g_1g_2\dots g_n = (\dots((g_1g_2)g_3)\dots)g_n.$
Любая другая расстановка скобок, например $(g_1g_2)(g_3g_4)$, в силу ассоциативности даёт тот же самый элемент группы.

\subsubsection*{Доказательство}

Доказательство проводится индукцией по $n$. При $n=2$ утверждение очевидно. Переход от $n-1$ к $n$ следует из аксиомы ассоциативности и предположения индукции, применённого к произведениям меньшей длины.

\subsection{Целая степень элемента группы}

Пусть $g \in G$ и $n \in \ZZ$.

\begin{Definition}{Степень элемента}{}
	Степень элемента $g^n$ определяется следующим образом:
	\[
	g^n =
	\begin{cases}
		\underbrace{g \cdot g \cdot \dots \cdot g}_{n}, & n > 0, \\
		e, & n = 0, \\
		\underbrace{g^{-1} \cdot g^{-1} \cdot \dots \cdot g^{-1}}_{|n|}, & n < 0.
	\end{cases}
	\]
\end{Definition}

\begin{Proposition}{Свойства степеней}{}
	Для любых $k,l \in \ZZ$:
	\begin{enumerate}
		\item $g^k g^l = g^{k+l}$;
		\item $(g^k)^l = g^{kl}$.
	\end{enumerate}
\end{Proposition}

\subsubsection*{Доказательство одного случая}

Рассмотрим случай $k > 0$, $l < 0$ и $k+l > 0$. Пусть $l = -t$, где $t \in \NN$. Тогда $k > t$ и $g^kg^l = g^kg^{-t}.$
После последовательного сокращения $t$ пар $g$ и $g^{-1}$ получаем $g^kg^{-t} = g^{k-t} = g^{k+l}.$
Остальные случаи рассматриваются аналогично.

\subsection{Порядок элемента и циклические подгруппы}

\begin{Definition}{Порядок элемента}{}
	Если существует $n \in \NN$ такое, что $g^n = e$, то говорят, что элемент $g$ имеет конечный порядок. Порядком элемента $g$ называется число $\operatorname{ord} g = \min\{n \in \NN \mid g^n = e\}.$
	Если такого $n$ не существует, то говорят, что $\operatorname{ord} g = \infty.$
\end{Definition}

\begin{Proposition}{}{}
	Пусть $\operatorname{ord} g = d < \infty$. Тогда $\{m \in \ZZ \mid g^m = e\} = d\ZZ.$
	То есть $g^m = e \iff d \mid m.$
\end{Proposition}

\subsubsection*{Доказательство}

Разделим $m$ на $d$ с остатком: $m = dq + r, \qquad 0 \le r < d.$
Тогда $g^m = e \implies g^{dq+r} = e \implies (g^d)^qg^r = e.$
Так как $g^d = e$, получаем $g^r = e$. Поскольку $d$ — минимальная натуральная степень, дающая $e$, а $0 \le r < d$, единственный возможный вариант — $r = 0$. Следовательно, $m = dq$, то есть $d \mid m$.

\subsection{Подгруппа, порождённая элементом}

Для любого элемента $g \in G$ можно рассмотреть множество всех его целых степеней.

\begin{Definition}{Циклическая подгруппа}{}
	Множество $\langle g \rangle = \{g^n \mid n \in \ZZ\}$
	называется циклической подгруппой, порождённой элементом $g$.
\end{Definition}

\subsubsection*{Проверка подгрупповых аксиом}

Проверим, что $\langle g \rangle$ действительно является подгруппой в $G$:
\begin{enumerate}
	\item $e \in \langle g \rangle$, так как $e = g^0$.
	
	\item Если $h \in \langle g \rangle$, то $h = g^n$ для некоторого $n \in \ZZ$. Тогда $h^{-1} = (g^n)^{-1} = g^{-n} \in \langle g \rangle.$
	
	\item Если $h_1,h_2 \in \langle g \rangle$, то $h_1 = g^{n_1}$ и $h_2 = g^{n_2}$ для некоторых $n_1,n_2 \in \ZZ$. Тогда $h_1h_2 = g^{n_1}g^{n_2} = g^{n_1+n_2} \in \langle g \rangle.$
\end{enumerate}

Таким образом, $\langle g \rangle < G.$

\begin{Proposition}{}{}
	Пусть $g \in G$.
	\begin{enumerate}
		\item Если $\operatorname{ord} g = \infty$, то $|\langle g \rangle| = \infty$ и все степени $g^i$ различны.
		\item Если $\operatorname{ord} g = n < \infty$, то $|\langle g \rangle| = n$
		и $\langle g \rangle = \{g^0,g^1,\dots,g^{n-1}\}.$
	\end{enumerate}
\end{Proposition}

\subsubsection*{Доказательство}

\begin{enumerate}
	\item Пусть $\operatorname{ord} g = \infty$. Предположим, что $g^i = g^j$ при $i > j$. Тогда $g^{i-j} = e.$
	Так как $i-j > 0$, это противоречит условию $\operatorname{ord} g = \infty$. Значит, все степени различны.
	
	\item Пусть $\operatorname{ord} g = n < \infty$. Для любого $m \in \ZZ$ представим $m$ в виде $m = nq + r, \qquad 0 \le r < n.$
	Тогда $g^m = (g^n)^qg^r = g^r.$
	Следовательно, $\langle g \rangle = \{g^0,g^1,\dots,g^{n-1}\}.$
	
	Покажем, что эти элементы различны. Если $g^i = g^j$ для $0 \le j < i \le n-1$, то $g^{i-j} = e.$
	Но $0 < i-j \le n-1$, что противоречит минимальности $n$.
\end{enumerate}

Таким образом, $|\langle g \rangle| = \operatorname{ord} g.$

\begin{Definition}{Циклическая группа}{}
	Если существует такой элемент $g \in G$, что $\langle g \rangle = G,$
	то группа $G$ называется циклической, а элемент $g$ — её образующим, или порождающим, элементом.
\end{Definition}

\subsubsection*{Таблица Кэли для циклической группы}

Для группы $\langle g \rangle = \{g^0,g^1,\dots,g^{n-1}\}$
умножение задаётся правилом $g^i \cdot g^j = g^{\,i+j \pmod n}.$
Поэтому таблица Кэли имеет сдвинутую структуру:
\[
\begin{pmatrix}
	g^0 & g^1 & g^2 & \dots & g^{n-1} \\
	g^1 & g^2 & g^3 & \dots & g^0 \\
	g^2 & g^3 & g^4 & \dots & g^1 \\
	\vdots & \vdots & \vdots & \ddots & \vdots \\
	g^{n-1} & g^0 & g^1 & \dots & g^{n-2}
\end{pmatrix}.
\]

\subsubsection*{Примеры}

Основной пример циклической группы порядка $n$ — группа вычетов по модулю $n$ по сложению: $(\ZZ/n\ZZ, +).$
Здесь $\ZZ/n\ZZ = \langle \bar{1} \rangle.$

\subsection{Критерий цикличности конечной группы}

Пример: $\ZZ^\times = \{\pm 1\},$
где $\ZZ^\times$ — группа обратимых элементов кольца целых чисел по умножению.

\begin{Proposition}{}{}
	Пусть $|G| = n < \infty$. Тогда $G$ является циклической тогда и только тогда, когда в ней существует элемент порядка $n$.
\end{Proposition}

\subsubsection*{Доказательство}

\begin{itemize}
	\item[$\Rightarrow$)] Если $G$ циклическая, то по определению $G = \langle a \rangle$ для некоторого $a \in G$. Тогда $\operatorname{ord} a = |\langle a \rangle| = |G| = n.$
	
	\item[$\Leftarrow$)] Пусть существует $g \in G$ такой, что $\operatorname{ord} g = n$. Тогда $|\langle g \rangle| = n.$
	Так как $\langle g \rangle \subseteq G$ и их порядки совпадают, то $\langle g \rangle = G.$
	Следовательно, $G$ циклическая.
\end{itemize}
